% Source: https://artofproblemsolving.com/community/c3h1565384_repeated_bags_of_marbles
% Stable identity: aops-topic-1565384-q01
The bag always holds $3$ marbles (two out, two in), so the state is the number $r$ of red marbles.
\begin{itemize}
\item $r=2$ (RRG): the pair is RR with probability $\dfrac{1}{3}$, giving R, G, R, i.e.\ $r=1$; it is RG with probability $\dfrac{2}{3}$, giving R, R, R, i.e.\ $r=3$.
\item $r=3$ (RRR): the pair is RR, leaving R and adding R, G, so $r=2$.
\item $r=1$ (RGG): the pair is RG (probability $\dfrac{2}{3}$), leaving G and adding R, R, so $r=2$; or GG, leaving R and adding R, G, so $r=2$. Either way $r=2$.
\end{itemize}
So from $r=1$ or $r=3$ the bag always returns to $r=2$. Starting at $r=2$, after turns $2$ and $4$ the bag is certainly at $r=2$ again, and turn $5$ is just like turn $1$: it ends with $r=3$ with probability $\dfrac{2}{3}$.

Hence $\dfrac {p}{q}=\dfrac{2}{3}$ and $p+q=\boxed{5}$.
